% Supplement subsections typed below as "Section~S<n>.<m>"; build.sh checks each number against the built supplement.
% supplement-ref supp:iia=S1.1
% supplement-ref supp:borda=S1.8
% supplement-ref supp:learned-development=S6.2
\section{Introduction}\label{sec:intro}

Most neural networks store their inputs and their learned parameters as numbers, and they learn by following gradients. This paper asks whether a network can instead compute with rankings, ordered lists of items, and learn by counting votes. It also asks whether the signal of those votes can train an ordinary neural network as well. Codes built from orders appear in neuroscience and in similarity search (Section~\ref{sec:related}). Here the network learns the orders that it compares its inputs with.

ArrowFlow builds such a network from four ingredients:
\begin{itemize}[leftmargin=2em,topsep=2pt,itemsep=1pt]
  \item \textbf{Data points} are rankings. A ranking is a permutation of a fixed set of items, the vocabulary.
  \item \textbf{Learned filters} are rankings of the same vocabulary. Each \emph{ranking filter} is a permutation prototype, a learned ranking that an input is compared with.
  \item \textbf{Similarity} is Spearman's footrule distance. It counts the total number of positions by which the items of one ranking must move to match another \citep{diaconis1977footrule}.
  \item \textbf{Learning} reorders the items of a filter by counting votes for their positions. No derivative is computed.
\end{itemize}

An encoder turns a real-valued example into a ranking (Figure~\ref{fig:pipeline}). It expands, standardizes and projects the features, and then it sorts the projected scores. A \emph{ranking layer} holds a bank of ranking filters and outputs the bank in order of distance to its input, nearest filter first. That output is itself a ranking, so the next layer can read it. To predict, ArrowFlow compares an example's last hidden ranking with those of the training examples, and the nearest ones vote on the class. We call this step the nearest-neighbor readout. Several \emph{views}, each an encoder with its own trained network, then vote on the final class.

Training is the main contribution of this paper: a layer of ranking filters are adjusted based on motion signals that are coming from higher layers. It is similar to backpropagation, but it is on ranking filters. It is the first known method to train a hierarchical network through discrete signals that are not derivatives.
During training, the network gets one more layer, an output layer with one filter per class. When this layer assigns an example to the wrong class, the filter of the true class votes to move toward the example's last hidden ranking. The class filter and the example's ranking are both rankings of the same hidden filters. The signed differences between the positions that they give each hidden filter are called \emph{motions}, and the motions travel back down the network. If the class filter ranks a hidden filter earlier than the example's ranking does, that hidden filter votes to move toward its input. If the class filter ranks it later, the filter votes to move away. Nothing is differentiated along the way. Still, whenever the class filter votes, the motions are exactly half the gradient of the squared distance between the two rankings' positions, taken with respect to the positions of the hidden filters in the example's ranking. A hidden filter's vote then tends to move its position the way a gradient step would. So the vote stands in for the missing derivative of the sort, as a straight-through estimator does \citep{bengio2013estimating}. Such an estimator trains through a discrete step by replacing the step's derivative with a simple surrogate (Section~S1.8).

\begin{figure}[!htbp]
\centering
\resizebox{\textwidth}{!}{%
\begin{tikzpicture}[
    >=Stealth,
    box/.style={draw, rounded corners=2pt, minimum height=0.7cm, minimum width=1.4cm,
                font=\small, align=center, thick},
    databox/.style={box, fill=blue!8},
    procbox/.style={box, fill=orange!12},
    sortbox/.style={box, fill=green!12},
    votebox/.style={box, fill=red!10},
    arr/.style={->, thick},
    lbl/.style={font=\footnotesize, text=black!70},
]
% === Shared encoding steps (top row); every view expands and standardizes the same way ===
\node[databox] (input) at (0,0) {$x \in \R^d$};
\node[procbox, right=0.6 of input] (poly) {Polynomial\\[-1pt]Expansion};
\node[procbox, right=0.6 of poly] (std) {Standardization};
\draw[arr] (input) -- (poly);
\draw[arr] (poly) -- (std);
\node[lbl, above=0.15 of input] {$d$ feat.};
\node[lbl, above=0.15 of poly] {$d'$ feat.};
% === One projection and argsort per view: view k ranks the standardized features with its own W_k ===
\node[procbox] (proj1) at (0,-1.75) {Projection $W_1$,\\[-1pt]$\argsort$};
\node[procbox] (proj2) at (4,-1.75) {Projection $W_2$,\\[-1pt]$\argsort$};
\node[procbox] (projK) at (8,-1.75) {Projection $W_K$,\\[-1pt]$\argsort$};
\node at ($(proj2)!0.5!(projK)$) {\Large$\cdots$};
\draw[arr] (std.south) -- ++(0,-0.45) -| (proj1.north);
\draw[arr] (std.south) -- ++(0,-0.45) -| (proj2.north);
\draw[arr] (std.south) -- ++(0,-0.45) -| (projK.north);
% === Multi-view branching (3 networks), each reading its own input ranking ===
\node[sortbox, minimum height=1.1cm] (net1) at (0, -3.6) {ArrowFlow\\[-1pt]Net 1};
\node[sortbox, minimum height=1.1cm] (net2) at (4, -3.6) {ArrowFlow\\[-1pt]Net 2};
\node[sortbox, minimum height=1.1cm] (netK) at (8, -3.6) {ArrowFlow\\[-1pt]Net $K$};
\node at ($(net2)!0.5!(netK)$) {\Large$\cdots$};
\draw[arr] (proj1) -- node[lbl, right] {$\pi_1 \in S_e$} (net1);
\draw[arr] (proj2) -- node[lbl, right] {$\pi_2 \in S_e$} (net2);
\draw[arr] (projK) -- node[lbl, right] {$\pi_K \in S_e$} (netK);
% === Predictions ===
\node[databox, below=0.6 of net1] (pred1) {$\hat{y}_1$};
\node[databox, below=0.6 of net2] (pred2) {$\hat{y}_2$};
\node[databox, below=0.6 of netK] (predK) {$\hat{y}_K$};
\draw[arr] (net1) -- (pred1);
\draw[arr] (net2) -- (pred2);
\draw[arr] (netK) -- (predK);
% === Majority vote ===
\node[votebox, below=0.7 of pred2, minimum width=2.8cm] (vote) {Majority vote};
\draw[arr] (pred1.south) -- ++(0,-0.2) -| (vote.160);
\draw[arr] (pred2) -- (vote);
\draw[arr] (predK.south) -- ++(0,-0.2) -| (vote.20);
\node[databox, below=0.5 of vote, fill=red!15] (final) {$\hat{y}$};
\draw[arr] (vote) -- (final);
% === Network detail inset (far right, well separated) ===
\begin{scope}[shift={(13.8,-1.9)}]
\node[font=\small\bfseries, anchor=south] at (0,0.5) {Network Detail};
\draw[thick, rounded corners=3pt, fill=green!4] (-3.35,-4.85) rectangle (2.85,0);
\coordinate (insetleft) at (-3.35,-2.4);
\node[sortbox, minimum width=2.6cm, font=\footnotesize] (L1) at (0,-0.55) {Ranking layer 1\\[-1pt]$N_1$ filters};
\node[sortbox, minimum width=2.6cm, font=\footnotesize] (L2) at (0,-1.95) {Ranking layer 2\\[-1pt]$N_2$ filters};
\draw[arr] (L1) -- node[lbl, right, font=\scriptsize] {$\pi'{\in}S_{N_1}$} (L2);
\node[lbl, above=0.08 of L1, font=\scriptsize] {$\pi_k \in S_e$};
% below the last hidden layer: the output layer, used in training, and the nearest-neighbor readout, used to predict
\node[sortbox, minimum width=2.1cm, fill=green!20, font=\scriptsize] (Lout) at (-1.25,-3.55) {Output layer $L$\\[-1pt]$C$ filters};
\node[votebox, minimum width=2.1cm, font=\scriptsize] (RO) at (1.25,-3.55) {kNN readout\\[-1pt]stored rankings};
\draw[arr, dashed] ([xshift=-4mm]L2.south) -- ++(0,-0.4) -| (Lout.north);
\draw[arr] ([xshift=4mm]L2.south) -- ++(0,-0.4) -| (RO.north);
% training: the motions returned by the output votes train the last hidden layer (dashed, upward), which relays signed
% motions to the layer below (dashed, between the two hidden layers)
\draw[dashed, thick] (Lout.west) -- (-2.85,-3.55) -- (-2.85,-1.95);
\draw[arr, dashed] (-2.85,-1.95) -- (L2.west);
\node[lbl, font=\scriptsize, rotate=90, anchor=south] at (-2.85,-2.75) {returned motions};
\draw[arr, dashed] ([xshift=-7mm]L2.north) -- node[lbl, left, font=\scriptsize] {relay} ([xshift=-7mm]L1.south);
\node[lbl, font=\scriptsize, anchor=east] at (-1.33,-2.3) {$\pi''{\in}S_{N_2}$};
\node[lbl, below=0.08 of Lout, font=\scriptsize, align=center] {training:\\[-1pt]votes};
\node[lbl, below=0.08 of RO, font=\scriptsize, align=center] {prediction: vote of the\\[-1pt]nearest stored rankings};
\end{scope}
% Dashed connector from network box to inset
\draw[dashed, gray, thick] (netK.east) -- ++(0.8,0) |- (insetleft);
\end{tikzpicture}%
}
\caption{\textbf{ArrowFlow pipeline.} A real-valued input $x$ is encoded by polynomial expansion of degree $p$, standardization, a projection and argsort. The $K$ views use different projections $W_k$, $k = 1, \dots, K$, cycling through target-aware, random and calibrated strategies, so their networks receive different input rankings. Each network (detail, right) is a stack of ranking layers, and every layer ranks its ranking filters by Spearman's footrule distance to its input ranking. In training, an output layer holds one filter per class, and the filter of the example's class votes to move toward the last hidden ranking. The motions this vote returns train the last hidden layer, which relays signed motions to the preceding layer (dashed arrows). To predict, a nearest-neighbor (kNN) readout compares the last hidden ranking with the stored hidden rankings of the training examples. The $K$ predictions are combined by majority vote. $S_n$ denotes the set of rankings of $n$ items.}
\label{fig:pipeline}
\end{figure}


Every update of a ranking filter merges the rankings of a training batch into one, so every update is an election. Arrow's impossibility theorem \citep{arrow1951social,arrow1963social} says that no way of merging rankings has every property one might want. With three or more alternatives, no rule that turns every profile of rankings into one ranking satisfies three conditions at once: Pareto efficiency, independence of irrelevant alternatives and non-dictatorship. This is a lens to look at the update rule. It gives us a simple perspective on the context dependence of a filter's order of two items, which is essential for nonlinearity in the network.
Independence asks that the merged order of two items depend only on how the voters order those two items. Our update is a weighted Borda count, which orders items by their weighted sums of positions, and the filter's current order counts as one more voter. The rule is Pareto efficient for any positive weights. Under a mild condition on the weights it is also non-dictatorial, so it must break independence. The condition says that the votes of a batch weigh enough that they can move the filter. In ArrowFlow's layers, an update that holds the current order fixed is not Pareto efficient, so the theorem does not apply to it. It breaks independence too, by a direct construction, whenever its votes can move the filter (Section~S1.1). So a filter's order of two items can change when the training rankings move a third.

We use the theorem as a design lens, and it promises nothing about accuracy. It shows that this context dependence lives in the learning rule. And the voting reading names three stages of the network whose rules can be exchanged. The name joins \emph{Arrow}, for the vote inside every filter update, and \emph{Flow}, for the rankings that pass from layer to layer.

\paragraph{What this paper shows.} We did not build ArrowFlow to be the most accurate classifier, and it is not one. With all settings tuned on the training data alone, it trails the best tuned classical model on all seventeen tabular datasets we test (Table~\ref{tab:main}). Its average rank is fourth among the six tuned models, between random forest and gradient boosting, and fifth on the ten further datasets alone. What the paper shows instead is that votes, without computing any derivative, can train ranking layers, and that the motions they return can train an ordinary neural network:
\begin{enumerate}[leftmargin=2em,topsep=2pt,itemsep=1pt]
  \item \emph{Training helps.} The trained network is more accurate than an untrained ArrowFlow on 16 of the 17 datasets, although the untrained version comes within one percentage point on 12 of the 17 (Section~\ref{sec:training-controls}).
  \item \emph{The training signal carries information.} Suppose we scramble the motions but keep their sizes. The network then ends up less accurate than if its hidden filters had never moved. This happens on 16 of the 17 datasets when each motion reaches the wrong filter, and on 14 of the 17 when its direction is random. A post hoc test across the datasets, one defined after the results were known, supports both (Section~\ref{sec:motion-controls}).
  \item \emph{Motions also train a real-valued network.} Motions of the same kind can be turned into target values for the scores that a small encoder network sorts. They then train that network with a standard gradient-based optimizer, although the sort between the two is never differentiated. Other methods also train through an exact discrete step (Section~\ref{sec:related}). The trained encoder network is more accurate than the same network untrained on all 17 datasets, significantly on 11. Swapped into ArrowFlow for the fixed encoder, it raises the mean accuracy on 11 of the 17 datasets, significantly on balance-scale, and lowers it significantly on none (Section~\ref{sec:learned}). The conversion of motions into target scores and its settings were developed on single splits of twelve of the seventeen datasets (Section~S6.2).
  \item \emph{The contribution is the learning rule, and its limits are clear.} A support vector classifier with a fixed Kendall kernel, a standard similarity between rankings, is more accurate than ArrowFlow on 14 of the 17 datasets when both read rankings from the same encoder family. Read by that classifier, trained hidden rankings are no better than untrained ones. So training improves the neighborhoods of the nearest-neighbor readout, but the representation read by a strong fixed readout shows no gain. A second hidden layer did not help (Section~\ref{sec:controlled}).
\end{enumerate}
Sections~\ref{sec:related} to~\ref{sec:theory} cover related work, the network, the Arrow lens and the theory, and Section~\ref{sec:experiments} reports the experiments. Section~\ref{sec:discussion} discusses the results, and Section~\ref{sec:conclusion} concludes. The supplement, Sections~S1 to~S6, holds the proofs, the method details and protocol, the complete results, the ablations, the reproducibility records and the details of the trained encoder.
